\section{El problema del procesamiento del lenguaje natural}
\label{sec:problema_nlp}

El procesamiento del lenguaje natural (NLP, por sus siglas en inglés) tiene
como objetivo dotar a los sistemas computacionales de la capacidad de
comprender y generar lenguaje humano. A diferencia de otros dominios donde
las entradas son vectores numéricos bien definidos, el lenguaje presenta
propiedades estructurales que lo hacen particularmente difícil de modelar
con las arquitecturas revisadas en la sección anterior.

Esta sección caracteriza las cuatro propiedades del lenguaje natural que
resultan más problemáticas para los modelos clásicos y que motivan
directamente el diseño del mecanismo de atención.

%------------------------------------------------------------
\subsection{El lenguaje como secuencia de símbolos discretos}
\label{subsec:lenguaje_secuencia}

El lenguaje natural se expresa como secuencias de unidades discretas:
palabras, subpalabras o caracteres, denominados genéricamente
\textit{tokens}. A diferencia de una señal continua, cada token pertenece
a un vocabulario finito $\mathcal{V}$ y su representación numérica inmediata
(\textit{one-hot}) no captura ninguna relación con los demás tokens.

Más aún, el lenguaje es fundamentalmente \textit{variable en longitud}: una
oración puede tener 3 tokens o 300. Esto contrasta con los MLPs, que
requieren entradas de dimensión fija, y con las RNNs, que procesan la
secuencia de forma estrictamente secuencial. Cualquier arquitectura que
pretenda modelar lenguaje debe ser capaz de operar sobre secuencias de
longitud arbitraria sin perder información posicional ni contextual.

%------------------------------------------------------------
\subsection{Ambigüedad semántica}
\label{subsec:ambiguedad}

Una de las propiedades más características del lenguaje natural es su
ambigüedad: una misma forma superficial puede tener significados
completamente distintos dependiendo del contexto en que aparece. Esta
ambigüedad ocurre a múltiples niveles:

\begin{itemize}
    \item \textbf{Léxica:} una misma palabra tiene distintos significados
    (\textit{polisemia}). Ejemplo: \textit{banco} puede ser una institución
    financiera o un asiento exterior.

    \item \textbf{Sintáctica:} una misma oración admite distintas
    estructuras gramaticales. Ejemplo: \textit{``Vi al espía con el
    telescopio''} puede significar que yo usé el telescopio, o que el
    espía lo tenía.

    \item \textbf{Pragmática:} el significado depende del contexto
    situacional más allá del texto mismo.
\end{itemize}

% ============================================================
% fig_ambiguedad.tex: Ambigüedad semántica según contexto
% Uso: % ============================================================
% fig_ambiguedad.tex: Ambigüedad semántica según contexto
% Uso: % ============================================================
% fig_ambiguedad.tex: Ambigüedad semántica según contexto
% Uso: \input{figures/fig_ambiguedad}
% ============================================================
\begin{figure}[H]
\centering
\begin{tikzpicture}[
    word/.style={draw=gray!40, rounded corners, fill=gray!8,
                 minimum width=1.1cm, minimum height=0.65cm, font=\small},
    key/.style={draw=blue!60, rounded corners, fill=blue!12,
                minimum width=1.1cm, minimum height=0.65cm, font=\small\bfseries},
    meaning/.style={draw, rounded corners, fill=white,
                    minimum width=2.6cm, minimum height=0.7cm,
                    font=\footnotesize, align=center},
    arrow/.style={-{Stealth}, thick},
    ctx/.style={font=\footnotesize\itshape, gray}
]

% ── Oración 1 ────────────────────────────────────────────────
\node[ctx] at (3.5, 5.6) {Contexto A:};
\foreach \w/\i in {Fui/0, al/1, banco/2, a/3, depositar/4, dinero/5}{
    \pgfmathparse{int(\i==2)}
    \ifnum\pgfmathresult=1
        \node[key]  at (\i*1.18, 4.95) {\w};
    \else
        \node[word] at (\i*1.18, 4.95) {\w};
    \fi
}
\draw[arrow, blue!60] (2.36, 4.62) -- (2.36, 4.0);
\node[meaning, fill=green!10, draw=green!50!black]
    at (2.36, 3.6) {institución\\financiera};

% ── Oración 2 ────────────────────────────────────────────────
\node[ctx] at (3.5, 2.9) {Contexto B:};
\foreach \w/\i in {Me/0, senté/1, en/2, el/3, banco/4, del/5, parque/6}{
    \pgfmathparse{int(\i==4)}
    \ifnum\pgfmathresult=1
        \node[key]  at (\i*1.1, 2.3) {\w};
    \else
        \node[word] at (\i*1.1, 2.3) {\w};
    \fi
}
\draw[arrow, blue!60] (4.4, 1.97) -- (4.4, 1.35);
\node[meaning, fill=orange!10, draw=orange!60]
    at (4.4, 0.95) {asiento\\exterior};

% ── Anotación central ────────────────────────────────────────
\node[draw=blue!20, fill=blue!5, rounded corners, font=\footnotesize,
      inner sep=6pt, align=center] at (7.8, 2.8)
    {\textbf{``banco''} es la misma\\palabra en ambos casos.\\El significado\\
     solo emerge del\\contexto circundante.};

\draw[dashed, gray!50] (5.9, 3.6) -- (6.8, 3.2);
\draw[dashed, gray!50] (5.8, 0.95) -- (6.8, 2.4);

\end{tikzpicture}
\caption{Ambigüedad semántica léxica. La palabra \textit{banco} tiene
significados completamente distintos según el contexto oracional. Un modelo
que procese cada token de forma independiente, sin considerar las palabras
circundantes, no puede resolver esta ambigüedad.}
\label{fig:ambiguedad}
\end{figure}

% ============================================================
\begin{figure}[H]
\centering
\begin{tikzpicture}[
    word/.style={draw=gray!40, rounded corners, fill=gray!8,
                 minimum width=1.1cm, minimum height=0.65cm, font=\small},
    key/.style={draw=blue!60, rounded corners, fill=blue!12,
                minimum width=1.1cm, minimum height=0.65cm, font=\small\bfseries},
    meaning/.style={draw, rounded corners, fill=white,
                    minimum width=2.6cm, minimum height=0.7cm,
                    font=\footnotesize, align=center},
    arrow/.style={-{Stealth}, thick},
    ctx/.style={font=\footnotesize\itshape, gray}
]

% ── Oración 1 ────────────────────────────────────────────────
\node[ctx] at (3.5, 5.6) {Contexto A:};
\foreach \w/\i in {Fui/0, al/1, banco/2, a/3, depositar/4, dinero/5}{
    \pgfmathparse{int(\i==2)}
    \ifnum\pgfmathresult=1
        \node[key]  at (\i*1.18, 4.95) {\w};
    \else
        \node[word] at (\i*1.18, 4.95) {\w};
    \fi
}
\draw[arrow, blue!60] (2.36, 4.62) -- (2.36, 4.0);
\node[meaning, fill=green!10, draw=green!50!black]
    at (2.36, 3.6) {institución\\financiera};

% ── Oración 2 ────────────────────────────────────────────────
\node[ctx] at (3.5, 2.9) {Contexto B:};
\foreach \w/\i in {Me/0, senté/1, en/2, el/3, banco/4, del/5, parque/6}{
    \pgfmathparse{int(\i==4)}
    \ifnum\pgfmathresult=1
        \node[key]  at (\i*1.1, 2.3) {\w};
    \else
        \node[word] at (\i*1.1, 2.3) {\w};
    \fi
}
\draw[arrow, blue!60] (4.4, 1.97) -- (4.4, 1.35);
\node[meaning, fill=orange!10, draw=orange!60]
    at (4.4, 0.95) {asiento\\exterior};

% ── Anotación central ────────────────────────────────────────
\node[draw=blue!20, fill=blue!5, rounded corners, font=\footnotesize,
      inner sep=6pt, align=center] at (7.8, 2.8)
    {\textbf{``banco''} es la misma\\palabra en ambos casos.\\El significado\\
     solo emerge del\\contexto circundante.};

\draw[dashed, gray!50] (5.9, 3.6) -- (6.8, 3.2);
\draw[dashed, gray!50] (5.8, 0.95) -- (6.8, 2.4);

\end{tikzpicture}
\caption{Ambigüedad semántica léxica. La palabra \textit{banco} tiene
significados completamente distintos según el contexto oracional. Un modelo
que procese cada token de forma independiente, sin considerar las palabras
circundantes, no puede resolver esta ambigüedad.}
\label{fig:ambiguedad}
\end{figure}

% ============================================================
\begin{figure}[H]
\centering
\begin{tikzpicture}[
    word/.style={draw=gray!40, rounded corners, fill=gray!8,
                 minimum width=1.1cm, minimum height=0.65cm, font=\small},
    key/.style={draw=blue!60, rounded corners, fill=blue!12,
                minimum width=1.1cm, minimum height=0.65cm, font=\small\bfseries},
    meaning/.style={draw, rounded corners, fill=white,
                    minimum width=2.6cm, minimum height=0.7cm,
                    font=\footnotesize, align=center},
    arrow/.style={-{Stealth}, thick},
    ctx/.style={font=\footnotesize\itshape, gray}
]

% ── Oración 1 ────────────────────────────────────────────────
\node[ctx] at (3.5, 5.6) {Contexto A:};
\foreach \w/\i in {Fui/0, al/1, banco/2, a/3, depositar/4, dinero/5}{
    \pgfmathparse{int(\i==2)}
    \ifnum\pgfmathresult=1
        \node[key]  at (\i*1.18, 4.95) {\w};
    \else
        \node[word] at (\i*1.18, 4.95) {\w};
    \fi
}
\draw[arrow, blue!60] (2.36, 4.62) -- (2.36, 4.0);
\node[meaning, fill=green!10, draw=green!50!black]
    at (2.36, 3.6) {institución\\financiera};

% ── Oración 2 ────────────────────────────────────────────────
\node[ctx] at (3.5, 2.9) {Contexto B:};
\foreach \w/\i in {Me/0, senté/1, en/2, el/3, banco/4, del/5, parque/6}{
    \pgfmathparse{int(\i==4)}
    \ifnum\pgfmathresult=1
        \node[key]  at (\i*1.1, 2.3) {\w};
    \else
        \node[word] at (\i*1.1, 2.3) {\w};
    \fi
}
\draw[arrow, blue!60] (4.4, 1.97) -- (4.4, 1.35);
\node[meaning, fill=orange!10, draw=orange!60]
    at (4.4, 0.95) {asiento\\exterior};

% ── Anotación central ────────────────────────────────────────
\node[draw=blue!20, fill=blue!5, rounded corners, font=\footnotesize,
      inner sep=6pt, align=center] at (7.8, 2.8)
    {\textbf{``banco''} es la misma\\palabra en ambos casos.\\El significado\\
     solo emerge del\\contexto circundante.};

\draw[dashed, gray!50] (5.9, 3.6) -- (6.8, 3.2);
\draw[dashed, gray!50] (5.8, 0.95) -- (6.8, 2.4);

\end{tikzpicture}
\caption{Ambigüedad semántica léxica. La palabra \textit{banco} tiene
significados completamente distintos según el contexto oracional. Un modelo
que procese cada token de forma independiente, sin considerar las palabras
circundantes, no puede resolver esta ambigüedad.}
\label{fig:ambiguedad}
\end{figure}


La consecuencia directa para el modelado es que no existe una representación
fija correcta para un token: su significado es una función de su contexto.
Un modelo que asigne a cada token un vector de representación estático
(independiente del resto de la secuencia) no puede resolver esta ambigüedad.
Esto descarta, en principio, cualquier arquitectura sin mecanismo de
contextualización dinámica.

%------------------------------------------------------------
\subsection{Dependencias de largo alcance}
\label{subsec:dependencias_largas}

En el lenguaje natural, dos elementos de una oración pueden estar
relacionados gramaticalmente o semánticamente aunque estén separados por
muchos otros tokens. Estas relaciones se denominan \textit{dependencias de
largo alcance} y son ubicuas en el lenguaje real~\cite{bengio2003neural}.

Un ejemplo paradigmático es la concordancia sujeto-verbo en oraciones con
cláusulas subordinadas intercaladas:

\begin{quote}
\textit{``El científico que ganó el premio Nobel de física presentó su investigación.''}
\end{quote}

El verbo \textit{presentó} debe concordar en número y persona con el sujeto
\textit{científico}, aunque ambos estén separados por la cláusula relativa
\textit{``que ganó el premio Nobel de física''} (6 tokens). Un modelo con
ventana de contexto fija menor que esa distancia nunca puede aprender esta
relación correctamente.

% ============================================================
% fig_dependencias.tex — Dependencias largas y correferencia
% Uso: % ============================================================
% fig_dependencias.tex — Dependencias largas y correferencia
% Uso: % ============================================================
% fig_dependencias.tex — Dependencias largas y correferencia
% Uso: \input{figures/fig_dependencias}
% ============================================================
\begin{figure}[H]
\centering
\begin{tikzpicture}[
    word/.style={draw=gray!35, rounded corners, fill=gray!8,
                 minimum width=1.05cm, minimum height=0.62cm,
                 font=\scriptsize, inner sep=2pt},
    subj/.style={draw=blue!60, rounded corners, fill=blue!12,
                 minimum width=1.05cm, minimum height=0.62cm,
                 font=\scriptsize\bfseries, inner sep=2pt},
    verb/.style={draw=red!60, rounded corners, fill=red!10,
                 minimum width=1.05cm, minimum height=0.62cm,
                 font=\scriptsize\bfseries, inner sep=2pt},
    pron/.style={draw=orange!70, rounded corners, fill=orange!12,
                 minimum width=1.05cm, minimum height=0.62cm,
                 font=\scriptsize\bfseries, inner sep=2pt},
    lbl/.style={font=\tiny, align=center}
]

% ── Tokens colocados manualmente (12 tokens) ─────────────────
\node[subj] (t0)  at (0*1.08, 3.2) {El};
\node[subj] (t1)  at (1*1.08, 3.2) {cient.};
\node[word] (t2)  at (2*1.08, 3.2) {que};
\node[word] (t3)  at (3*1.08, 3.2) {gan\'o};
\node[word] (t4)  at (4*1.08, 3.2) {el};
\node[word] (t5)  at (5*1.08, 3.2) {premio};
\node[word] (t6)  at (6*1.08, 3.2) {Nobel};
\node[word] (t7)  at (7*1.08, 3.2) {de};
\node[word] (t8)  at (8*1.08, 3.2) {f\'isica};
\node[verb] (t9)  at (9*1.08, 3.2) {present\'o};
\node[pron] (t10) at (10*1.08,3.2) {su};
\node[word] (t11) at (11*1.08,3.2) {invest.};

% ── Arco: sujeto → verbo (dependencia larga) ─────────────────
\draw[-{Stealth}, blue!60, thick]
    (t1.north) .. controls (1.08, 4.6) and (9.72, 4.6) .. (t9.north);
\node[lbl, blue!70] at (5.4, 4.85)
    {dependencia sujeto--verbo \\ (8 tokens de distancia)};

% ── Arco: pronombre → sujeto (correferencia) ─────────────────
\draw[-{Stealth}, orange!80, thick]
    (t10.south) .. controls (10.8, 1.8) and (1.08, 1.8) .. (t1.south);
\node[lbl, orange!70] at (5.4, 1.4)
    {\textit{``su''} refiere a \textit{``cient.''} \\ (correferencia a distancia)};

% ── Leyenda ──────────────────────────────────────────────────
\node[subj, label=right:{\scriptsize sujeto}]    at (0.0,  0.55) {};
\node[verb, label=right:{\scriptsize verbo}]     at (2.5,  0.55) {};
\node[pron, label=right:{\scriptsize pronombre}] at (5.0,  0.55) {};

% ── Nota ─────────────────────────────────────────────────────
\node[draw=gray!30, fill=gray!5, rounded corners, font=\tiny,
      inner sep=5pt, align=left] at (9.5, 0.55)
    {Un MLP con ventana fija no puede\\ conectar tokens tan alejados.\\ Un Transformer s\'i: la atenci\'on es global.};

\end{tikzpicture}
\caption{Dependencia de largo alcance y correferencia. El verbo
\textit{present\'o} concuerda con \textit{cient\'ifico}, su sujeto,
separados por 8 tokens. El pronombre \textit{su} refiere tambi\'en al mismo
sujeto. Ambas relaciones requieren acceder a contexto distante.}
\label{fig:dependencias}
\end{figure}

% ============================================================
\begin{figure}[H]
\centering
\begin{tikzpicture}[
    word/.style={draw=gray!35, rounded corners, fill=gray!8,
                 minimum width=1.05cm, minimum height=0.62cm,
                 font=\scriptsize, inner sep=2pt},
    subj/.style={draw=blue!60, rounded corners, fill=blue!12,
                 minimum width=1.05cm, minimum height=0.62cm,
                 font=\scriptsize\bfseries, inner sep=2pt},
    verb/.style={draw=red!60, rounded corners, fill=red!10,
                 minimum width=1.05cm, minimum height=0.62cm,
                 font=\scriptsize\bfseries, inner sep=2pt},
    pron/.style={draw=orange!70, rounded corners, fill=orange!12,
                 minimum width=1.05cm, minimum height=0.62cm,
                 font=\scriptsize\bfseries, inner sep=2pt},
    lbl/.style={font=\tiny, align=center}
]

% ── Tokens colocados manualmente (12 tokens) ─────────────────
\node[subj] (t0)  at (0*1.08, 3.2) {El};
\node[subj] (t1)  at (1*1.08, 3.2) {cient.};
\node[word] (t2)  at (2*1.08, 3.2) {que};
\node[word] (t3)  at (3*1.08, 3.2) {gan\'o};
\node[word] (t4)  at (4*1.08, 3.2) {el};
\node[word] (t5)  at (5*1.08, 3.2) {premio};
\node[word] (t6)  at (6*1.08, 3.2) {Nobel};
\node[word] (t7)  at (7*1.08, 3.2) {de};
\node[word] (t8)  at (8*1.08, 3.2) {f\'isica};
\node[verb] (t9)  at (9*1.08, 3.2) {present\'o};
\node[pron] (t10) at (10*1.08,3.2) {su};
\node[word] (t11) at (11*1.08,3.2) {invest.};

% ── Arco: sujeto → verbo (dependencia larga) ─────────────────
\draw[-{Stealth}, blue!60, thick]
    (t1.north) .. controls (1.08, 4.6) and (9.72, 4.6) .. (t9.north);
\node[lbl, blue!70] at (5.4, 4.85)
    {dependencia sujeto--verbo \\ (8 tokens de distancia)};

% ── Arco: pronombre → sujeto (correferencia) ─────────────────
\draw[-{Stealth}, orange!80, thick]
    (t10.south) .. controls (10.8, 1.8) and (1.08, 1.8) .. (t1.south);
\node[lbl, orange!70] at (5.4, 1.4)
    {\textit{``su''} refiere a \textit{``cient.''} \\ (correferencia a distancia)};

% ── Leyenda ──────────────────────────────────────────────────
\node[subj, label=right:{\scriptsize sujeto}]    at (0.0,  0.55) {};
\node[verb, label=right:{\scriptsize verbo}]     at (2.5,  0.55) {};
\node[pron, label=right:{\scriptsize pronombre}] at (5.0,  0.55) {};

% ── Nota ─────────────────────────────────────────────────────
\node[draw=gray!30, fill=gray!5, rounded corners, font=\tiny,
      inner sep=5pt, align=left] at (9.5, 0.55)
    {Un MLP con ventana fija no puede\\ conectar tokens tan alejados.\\ Un Transformer s\'i: la atenci\'on es global.};

\end{tikzpicture}
\caption{Dependencia de largo alcance y correferencia. El verbo
\textit{present\'o} concuerda con \textit{cient\'ifico}, su sujeto,
separados por 8 tokens. El pronombre \textit{su} refiere tambi\'en al mismo
sujeto. Ambas relaciones requieren acceder a contexto distante.}
\label{fig:dependencias}
\end{figure}

% ============================================================
\begin{figure}[H]
\centering
\begin{tikzpicture}[
    word/.style={draw=gray!35, rounded corners, fill=gray!8,
                 minimum width=1.05cm, minimum height=0.62cm,
                 font=\scriptsize, inner sep=2pt},
    subj/.style={draw=blue!60, rounded corners, fill=blue!12,
                 minimum width=1.05cm, minimum height=0.62cm,
                 font=\scriptsize\bfseries, inner sep=2pt},
    verb/.style={draw=red!60, rounded corners, fill=red!10,
                 minimum width=1.05cm, minimum height=0.62cm,
                 font=\scriptsize\bfseries, inner sep=2pt},
    pron/.style={draw=orange!70, rounded corners, fill=orange!12,
                 minimum width=1.05cm, minimum height=0.62cm,
                 font=\scriptsize\bfseries, inner sep=2pt},
    lbl/.style={font=\tiny, align=center}
]

% ── Tokens colocados manualmente (12 tokens) ─────────────────
\node[subj] (t0)  at (0*1.08, 3.2) {El};
\node[subj] (t1)  at (1*1.08, 3.2) {cient.};
\node[word] (t2)  at (2*1.08, 3.2) {que};
\node[word] (t3)  at (3*1.08, 3.2) {gan\'o};
\node[word] (t4)  at (4*1.08, 3.2) {el};
\node[word] (t5)  at (5*1.08, 3.2) {premio};
\node[word] (t6)  at (6*1.08, 3.2) {Nobel};
\node[word] (t7)  at (7*1.08, 3.2) {de};
\node[word] (t8)  at (8*1.08, 3.2) {f\'isica};
\node[verb] (t9)  at (9*1.08, 3.2) {present\'o};
\node[pron] (t10) at (10*1.08,3.2) {su};
\node[word] (t11) at (11*1.08,3.2) {invest.};

% ── Arco: sujeto → verbo (dependencia larga) ─────────────────
\draw[-{Stealth}, blue!60, thick]
    (t1.north) .. controls (1.08, 4.6) and (9.72, 4.6) .. (t9.north);
\node[lbl, blue!70] at (5.4, 4.85)
    {dependencia sujeto--verbo \\ (8 tokens de distancia)};

% ── Arco: pronombre → sujeto (correferencia) ─────────────────
\draw[-{Stealth}, orange!80, thick]
    (t10.south) .. controls (10.8, 1.8) and (1.08, 1.8) .. (t1.south);
\node[lbl, orange!70] at (5.4, 1.4)
    {\textit{``su''} refiere a \textit{``cient.''} \\ (correferencia a distancia)};

% ── Leyenda ──────────────────────────────────────────────────
\node[subj, label=right:{\scriptsize sujeto}]    at (0.0,  0.55) {};
\node[verb, label=right:{\scriptsize verbo}]     at (2.5,  0.55) {};
\node[pron, label=right:{\scriptsize pronombre}] at (5.0,  0.55) {};

% ── Nota ─────────────────────────────────────────────────────
\node[draw=gray!30, fill=gray!5, rounded corners, font=\tiny,
      inner sep=5pt, align=left] at (9.5, 0.55)
    {Un MLP con ventana fija no puede\\ conectar tokens tan alejados.\\ Un Transformer s\'i: la atenci\'on es global.};

\end{tikzpicture}
\caption{Dependencia de largo alcance y correferencia. El verbo
\textit{present\'o} concuerda con \textit{cient\'ifico}, su sujeto,
separados por 8 tokens. El pronombre \textit{su} refiere tambi\'en al mismo
sujeto. Ambas relaciones requieren acceder a contexto distante.}
\label{fig:dependencias}
\end{figure}


Formalmente, si definimos la \textit{longitud de dependencia} entre dos
tokens $i$ y $j$ como $|i - j|$, los modelos de ventana fija de ancho $k$
son ciegos a toda dependencia con $|i - j| > k$. En textos reales, las
dependencias relevantes pueden extenderse a decenas de tokens~\cite{vaswani2017attention}.

%------------------------------------------------------------
\subsection{Orden y estructura sintáctica}
\label{subsec:orden}

El orden de los tokens en una secuencia es portador de significado: no es
simplemente un atributo accidental sino una dimensión informativa
fundamental. Considérense los pares:

\begin{center}
\begin{tabular}{ll}
\textit{``El perro muerde al hombre''} & $\neq$ \textit{``El hombre muerde al perro''} \\[2pt]
\textit{``No es verdad que miente''}   & $\neq$ \textit{``Es verdad que no miente''}   \\
\end{tabular}
\end{center}

Ambos pares contienen exactamente los mismos tokens; el significado cambia
radicalmente con el orden. Esto implica que cualquier modelo que procese los
tokens como un conjunto desordenado (es decir, que sea \textit{invariante a
permutaciones}) es estructuralmente incapaz de modelar lenguaje de forma
correcta~\cite{goodfellow2016deeplearning}.

Los MLPs son invariantes a permutaciones salvo que la posición esté
codificada explícitamente en la entrada. Las RNNs incorporan el orden
implícitamente a través del estado oculto secuencial, pero al costo de
procesar los tokens uno a uno. Los Transformers deben resolver este problema
de una tercera forma: codificando el orden \textit{explícitamente} sin
sacrificar el paralelismo. Esta solución, la codificación posicional, se
describe en la Sección~\ref{sec:arquitectura_transformer}.

%------------------------------------------------------------
\subsection{Correferencialidad}
\label{subsec:correferencia}

La \textit{correferencialidad} ocurre cuando distintas expresiones en un
texto refieren a la misma entidad. El caso más común son los pronombres:

\begin{quote}
\textit{``La investigadora publicó sus resultados. \textbf{Ella} había
trabajado durante diez años en el problema.''}
\end{quote}

Para interpretar correctamente \textit{ella}, el modelo debe establecer que
este pronombre corefiere con \textit{investigadora}, mencionada en la oración
anterior. Esta relación puede extenderse a través de múltiples oraciones y
requiere rastrear entidades a lo largo de un contexto potencialmente largo.

La correferencia es conceptualmente análoga al problema de Hopfield: así
como en una red de Hopfield el estado de cada unidad está determinado por el
estado global del sistema, resolver la correferencia requiere que cada token
sea interpretado en función de todos los demás. La diferencia clave es la
escala y la dirección: mientras Hopfield opera sobre patrones binarios fijos,
el lenguaje requiere relaciones dinámicas y asimétricas entre tokens.

%------------------------------------------------------------
\subsection{Análisis integrado: por qué el contexto es indispensable}
\label{subsec:contexto_indispensable}

Las cuatro propiedades anteriores convergen en una conclusión unificada:
\textbf{el significado de un token no puede determinarse en aislamiento}. Es
siempre una función del contexto en que aparece.

La siguiente figura analiza paso a paso cómo el contexto es imprescindible
para resolver incluso una ambigüedad simple:

% ============================================================
% fig_ejemplo_contexto.tex — Análisis paso a paso de una oración
% Uso: % ============================================================
% fig_ejemplo_contexto.tex — Análisis paso a paso de una oración
% Uso: % ============================================================
% fig_ejemplo_contexto.tex — Análisis paso a paso de una oración
% Uso: \input{figures/fig_ejemplo_contexto}
% ============================================================
\begin{figure}[H]
\centering
\begin{tikzpicture}[
    word/.style={draw=gray!40, rounded corners, fill=gray!8,
                 minimum width=1.0cm, minimum height=0.6cm, font=\small},
    focus/.style={draw=blue!60, rounded corners, fill=blue!12,
                  minimum width=1.0cm, minimum height=0.6cm,
                  font=\small\bfseries},
    ctx/.style={draw=green!50!black, rounded corners, fill=green!8,
                minimum width=1.0cm, minimum height=0.6cm, font=\small},
    box/.style={draw=gray!40, rounded corners, fill=white,
                inner sep=6pt, font=\footnotesize, align=left},
    arrow/.style={-{Stealth}, thick, gray},
    step/.style={font=\small\bfseries, blue!70}
]

% ── Oración completa (referencia) ───────────────────────────
\node[font=\small\bfseries, gray] at (3.5, 6.8)
    {Oración: \textit{``Vi a la estrella con el telescopio''}};

% ── PASO 1: token aislado ────────────────────────────────────
\node[step] at (-0.5, 5.9) {Paso 1};
\node[box]  at (4.2, 5.85) {Sin contexto: \textbf{``estrella''} puede ser\\
    \textcolor{blue!70}{$\bullet$ astro celeste} \quad
    \textcolor{orange!80}{$\bullet$ persona famosa}};

\node[focus] at (1.8, 5.85) {estrella};
\node[font=\tiny, red!60] at (1.8, 5.45) {ambigua};

% ── PASO 2: ventana pequeña (±2) ─────────────────────────────
\node[step] at (-0.5, 4.6) {Paso 2};
\node[box]  at (4.5, 4.55) {Ventana $\pm 2$: \textit{``la estrella con el''}\\
    Sigue ambigua: sin información resolutoria.};

\foreach \w/\i/\st in {la/0/ctx, estrella/1/focus, con/2/ctx, el/3/ctx}{
    \pgfmathparse{int(\i==1)}
    \ifnum\pgfmathresult=1
        \node[focus] at (\i*1.1, 4.55) {\w};
    \else
        \node[ctx]   at (\i*1.1, 4.55) {\w};
    \fi
}
\node[font=\tiny, red!60] at (1.1, 4.15) {aún ambigua};

% ── PASO 3: oración completa ────────────────────────────────
\node[step] at (-0.5, 3.2) {Paso 3};
\node[box, fill=green!5] at (5.0, 3.15)
    {Oración completa: \textit{``con el telescopio''}\\
    $\Rightarrow$ instrumento óptico $\Rightarrow$
    \textcolor{green!50!black}{\textbf{astro celeste}}};

\foreach \w/\i in {Vi/0, a/1, la/2, estrella/3, con/4, el/5, telescopio/6}{
    \pgfmathparse{int(\i==3)}
    \ifnum\pgfmathresult=1
        \node[focus] at (\i*1.05, 3.15) {\w};
    \else
        \pgfmathparse{int(\i==6 || \i==5 || \i==4)}
        \ifnum\pgfmathresult=1
            \node[ctx, draw=green!50!black, fill=green!10] at (\i*1.05, 3.15) {\w};
        \else
            \node[word] at (\i*1.05, 3.15) {\w};
        \fi
    \fi
}

% Arco resolutorio
\draw[-{Stealth}, green!60!black, thick]
    (6.3, 2.82) .. controls (6.3, 2.2) and (3.15, 2.2) .. (3.15, 2.82);
\node[font=\tiny, green!50!black] at (4.7, 1.95)
    {\textit{``telescopio''} resuelve la ambigüedad de \textit{``estrella''}};

% ── Conclusión ───────────────────────────────────────────────
\node[draw=blue!25, fill=blue!5, rounded corners, font=\footnotesize,
      inner sep=7pt, align=center] at (3.5, 1.1)
    {El significado correcto requiere acceso a \textbf{toda} la oración,\\
     no solo a una ventana local. Un modelo que no tenga acceso\\
     global al contexto no puede resolver esta ambigüedad.};

\end{tikzpicture}
\caption{Análisis paso a paso de la resolución de ambigüedad en la oración
\textit{``Vi a la estrella con el telescopio''}. En el Paso 1, el token
\textit{estrella} es ambiguo en aislamiento. En el Paso 2, una ventana local
no aporta información resolutoria. Solo en el Paso 3, con acceso a toda la
oración, la palabra \textit{telescopio} permite inferir que \textit{estrella}
se refiere a un astro y no a una persona famosa.}
\label{fig:ejemplo_contexto}
\end{figure}

% ============================================================
\begin{figure}[H]
\centering
\begin{tikzpicture}[
    word/.style={draw=gray!40, rounded corners, fill=gray!8,
                 minimum width=1.0cm, minimum height=0.6cm, font=\small},
    focus/.style={draw=blue!60, rounded corners, fill=blue!12,
                  minimum width=1.0cm, minimum height=0.6cm,
                  font=\small\bfseries},
    ctx/.style={draw=green!50!black, rounded corners, fill=green!8,
                minimum width=1.0cm, minimum height=0.6cm, font=\small},
    box/.style={draw=gray!40, rounded corners, fill=white,
                inner sep=6pt, font=\footnotesize, align=left},
    arrow/.style={-{Stealth}, thick, gray},
    step/.style={font=\small\bfseries, blue!70}
]

% ── Oración completa (referencia) ───────────────────────────
\node[font=\small\bfseries, gray] at (3.5, 6.8)
    {Oración: \textit{``Vi a la estrella con el telescopio''}};

% ── PASO 1: token aislado ────────────────────────────────────
\node[step] at (-0.5, 5.9) {Paso 1};
\node[box]  at (4.2, 5.85) {Sin contexto: \textbf{``estrella''} puede ser\\
    \textcolor{blue!70}{$\bullet$ astro celeste} \quad
    \textcolor{orange!80}{$\bullet$ persona famosa}};

\node[focus] at (1.8, 5.85) {estrella};
\node[font=\tiny, red!60] at (1.8, 5.45) {ambigua};

% ── PASO 2: ventana pequeña (±2) ─────────────────────────────
\node[step] at (-0.5, 4.6) {Paso 2};
\node[box]  at (4.5, 4.55) {Ventana $\pm 2$: \textit{``la estrella con el''}\\
    Sigue ambigua: sin información resolutoria.};

\foreach \w/\i/\st in {la/0/ctx, estrella/1/focus, con/2/ctx, el/3/ctx}{
    \pgfmathparse{int(\i==1)}
    \ifnum\pgfmathresult=1
        \node[focus] at (\i*1.1, 4.55) {\w};
    \else
        \node[ctx]   at (\i*1.1, 4.55) {\w};
    \fi
}
\node[font=\tiny, red!60] at (1.1, 4.15) {aún ambigua};

% ── PASO 3: oración completa ────────────────────────────────
\node[step] at (-0.5, 3.2) {Paso 3};
\node[box, fill=green!5] at (5.0, 3.15)
    {Oración completa: \textit{``con el telescopio''}\\
    $\Rightarrow$ instrumento óptico $\Rightarrow$
    \textcolor{green!50!black}{\textbf{astro celeste}}};

\foreach \w/\i in {Vi/0, a/1, la/2, estrella/3, con/4, el/5, telescopio/6}{
    \pgfmathparse{int(\i==3)}
    \ifnum\pgfmathresult=1
        \node[focus] at (\i*1.05, 3.15) {\w};
    \else
        \pgfmathparse{int(\i==6 || \i==5 || \i==4)}
        \ifnum\pgfmathresult=1
            \node[ctx, draw=green!50!black, fill=green!10] at (\i*1.05, 3.15) {\w};
        \else
            \node[word] at (\i*1.05, 3.15) {\w};
        \fi
    \fi
}

% Arco resolutorio
\draw[-{Stealth}, green!60!black, thick]
    (6.3, 2.82) .. controls (6.3, 2.2) and (3.15, 2.2) .. (3.15, 2.82);
\node[font=\tiny, green!50!black] at (4.7, 1.95)
    {\textit{``telescopio''} resuelve la ambigüedad de \textit{``estrella''}};

% ── Conclusión ───────────────────────────────────────────────
\node[draw=blue!25, fill=blue!5, rounded corners, font=\footnotesize,
      inner sep=7pt, align=center] at (3.5, 1.1)
    {El significado correcto requiere acceso a \textbf{toda} la oración,\\
     no solo a una ventana local. Un modelo que no tenga acceso\\
     global al contexto no puede resolver esta ambigüedad.};

\end{tikzpicture}
\caption{Análisis paso a paso de la resolución de ambigüedad en la oración
\textit{``Vi a la estrella con el telescopio''}. En el Paso 1, el token
\textit{estrella} es ambiguo en aislamiento. En el Paso 2, una ventana local
no aporta información resolutoria. Solo en el Paso 3, con acceso a toda la
oración, la palabra \textit{telescopio} permite inferir que \textit{estrella}
se refiere a un astro y no a una persona famosa.}
\label{fig:ejemplo_contexto}
\end{figure}

% ============================================================
\begin{figure}[H]
\centering
\begin{tikzpicture}[
    word/.style={draw=gray!40, rounded corners, fill=gray!8,
                 minimum width=1.0cm, minimum height=0.6cm, font=\small},
    focus/.style={draw=blue!60, rounded corners, fill=blue!12,
                  minimum width=1.0cm, minimum height=0.6cm,
                  font=\small\bfseries},
    ctx/.style={draw=green!50!black, rounded corners, fill=green!8,
                minimum width=1.0cm, minimum height=0.6cm, font=\small},
    box/.style={draw=gray!40, rounded corners, fill=white,
                inner sep=6pt, font=\footnotesize, align=left},
    arrow/.style={-{Stealth}, thick, gray},
    step/.style={font=\small\bfseries, blue!70}
]

% ── Oración completa (referencia) ───────────────────────────
\node[font=\small\bfseries, gray] at (3.5, 6.8)
    {Oración: \textit{``Vi a la estrella con el telescopio''}};

% ── PASO 1: token aislado ────────────────────────────────────
\node[step] at (-0.5, 5.9) {Paso 1};
\node[box]  at (4.2, 5.85) {Sin contexto: \textbf{``estrella''} puede ser\\
    \textcolor{blue!70}{$\bullet$ astro celeste} \quad
    \textcolor{orange!80}{$\bullet$ persona famosa}};

\node[focus] at (1.8, 5.85) {estrella};
\node[font=\tiny, red!60] at (1.8, 5.45) {ambigua};

% ── PASO 2: ventana pequeña (±2) ─────────────────────────────
\node[step] at (-0.5, 4.6) {Paso 2};
\node[box]  at (4.5, 4.55) {Ventana $\pm 2$: \textit{``la estrella con el''}\\
    Sigue ambigua: sin información resolutoria.};

\foreach \w/\i/\st in {la/0/ctx, estrella/1/focus, con/2/ctx, el/3/ctx}{
    \pgfmathparse{int(\i==1)}
    \ifnum\pgfmathresult=1
        \node[focus] at (\i*1.1, 4.55) {\w};
    \else
        \node[ctx]   at (\i*1.1, 4.55) {\w};
    \fi
}
\node[font=\tiny, red!60] at (1.1, 4.15) {aún ambigua};

% ── PASO 3: oración completa ────────────────────────────────
\node[step] at (-0.5, 3.2) {Paso 3};
\node[box, fill=green!5] at (5.0, 3.15)
    {Oración completa: \textit{``con el telescopio''}\\
    $\Rightarrow$ instrumento óptico $\Rightarrow$
    \textcolor{green!50!black}{\textbf{astro celeste}}};

\foreach \w/\i in {Vi/0, a/1, la/2, estrella/3, con/4, el/5, telescopio/6}{
    \pgfmathparse{int(\i==3)}
    \ifnum\pgfmathresult=1
        \node[focus] at (\i*1.05, 3.15) {\w};
    \else
        \pgfmathparse{int(\i==6 || \i==5 || \i==4)}
        \ifnum\pgfmathresult=1
            \node[ctx, draw=green!50!black, fill=green!10] at (\i*1.05, 3.15) {\w};
        \else
            \node[word] at (\i*1.05, 3.15) {\w};
        \fi
    \fi
}

% Arco resolutorio
\draw[-{Stealth}, green!60!black, thick]
    (6.3, 2.82) .. controls (6.3, 2.2) and (3.15, 2.2) .. (3.15, 2.82);
\node[font=\tiny, green!50!black] at (4.7, 1.95)
    {\textit{``telescopio''} resuelve la ambigüedad de \textit{``estrella''}};

% ── Conclusión ───────────────────────────────────────────────
\node[draw=blue!25, fill=blue!5, rounded corners, font=\footnotesize,
      inner sep=7pt, align=center] at (3.5, 1.1)
    {El significado correcto requiere acceso a \textbf{toda} la oración,\\
     no solo a una ventana local. Un modelo que no tenga acceso\\
     global al contexto no puede resolver esta ambigüedad.};

\end{tikzpicture}
\caption{Análisis paso a paso de la resolución de ambigüedad en la oración
\textit{``Vi a la estrella con el telescopio''}. En el Paso 1, el token
\textit{estrella} es ambiguo en aislamiento. En el Paso 2, una ventana local
no aporta información resolutoria. Solo en el Paso 3, con acceso a toda la
oración, la palabra \textit{telescopio} permite inferir que \textit{estrella}
se refiere a un astro y no a una persona famosa.}
\label{fig:ejemplo_contexto}
\end{figure}


Este ejemplo ilustra un principio general: un modelo que procese cada token
de forma independiente, o que solo tenga acceso a una ventana local, no puede
garantizar la interpretación correcta. Se requiere un mecanismo que permita
que \textit{cualquier} token acceda a la información de \textit{cualquier
otro} token de la secuencia, sin restricción de distancia.

%------------------------------------------------------------
\subsection{Requerimientos para un modelo de lenguaje adecuado}
\label{subsec:requerimientos}

Las propiedades del lenguaje descritas en esta sección imponen requerimientos
concretos a cualquier arquitectura que pretenda modelarlo correctamente.
La Tabla~\ref{tab:requerimientos_nlp} los resume y los contrasta con las
capacidades de los modelos clásicos ya analizados.

\begin{table}[H]
\centering
\caption{Requerimientos del modelado del lenguaje natural y cobertura por
arquitectura.}
\label{tab:requerimientos_nlp}
\begin{tabularx}{\textwidth}{lXcc}
\toprule
\textbf{Requerimiento} & \textbf{Motivación} & \textbf{Modelos clásicos} & \textbf{Transformer} \\
\midrule
Longitud variable        & Secuencias arbitrarias        & Limitado / No & Sí \\
Representación contextual & Ambigüedad semántica          & No            & Sí \\
Dependencias globales    & Largo alcance, correferencia  & No            & Sí \\
Sensibilidad al orden    & Estructura sintáctica         & Parcial       & Sí \\
Paralelización           & Eficiencia en entrenamiento   & No (RNN)      & Sí \\
\bottomrule
\end{tabularx}
\end{table}

Ninguna de las arquitecturas clásicas cubre simultáneamente todos estos
requerimientos. La siguiente sección introduce el mecanismo que permite
satisfacerlos: la \textit{atención}.
